\chapter{Perverse Degree and Weight Bookkeeping}
\label{chap:perverse-weights}

This chapter translates the punctual weights of
\cref{chap:weights} into perverse degree.  It proves only the formal
translation.  Stability of weights under \(Rf_!\), purity of intermediate
extensions, and the decomposition theorem are deliberately not used here;
they are downstream consequences of the Deligne--Laumon estimate.

Fix \(\iota:\Qbarl\simeq\mathbf C\).  All weight assertions in this chapter
are \(\iota\)-weight assertions unless “for every \(\iota\)” is stated.

\section{Weights of complexes}

\begin{definition}[Upper and lower weights]
\label{def:derived-weight-bounds}
A mixed Weil complex \(\cK_0\in\Dbc(X_0,\Qbarl)\) has weights at most \(w\)
if
\[
  \cH^a(\cK_0)\quad\text{has weights at most }w+a
  \quad\text{for every }a.
\tag{10.1}\label{eq:ordinary-derived-upper}
\]
Write \(\cK_0\in D_{\leq w}(X_0)\).  It has weights at least \(w\) if
\[
  \D_{X_0}\cK_0\in D_{\leq-w}(X_0),
\tag{10.2}\label{eq:derived-lower-dual}
\]
and is pure of weight \(w\) if it belongs to both
\(D_{\leq w}\) and \(D_{\geq w}\).
\end{definition}

The definition is engineered so that
\[
 \cK[n]\text{ pure of weight }w+n,\qquad
 \cK(m)\text{ pure of weight }w-2m,\qquad
 \D_X\cK\text{ pure of weight }-w.
\tag{10.3}\label{eq:weight-shift-twist-dual}
\]
Indeed, \(\cH^a(\cK[n])=\cH^{a+n}(\cK)\), geometric Frobenius on
\(\Qbarl(m)\) has weight \(-2m\), and the last assertion is built into
\eqref{eq:derived-lower-dual}.

\begin{warning}
“Punctually pure” is a property of a sheaf and its stalk eigenvalues.
“Pure complex” is the simultaneous upper/lower condition
\cref{def:derived-weight-bounds}.  They agree for a lisse sheaf in ordinary
degree zero on a smooth stratum, but they are not interchangeable words on a
singular space.
\end{warning}

\section{The dimension shift once more}

\begin{proposition}[Perverse normalization of a local system]
\label{prop:lisse-perverse-weight-shift}
Let \(U_0\) be smooth of pure dimension \(d\), and let \(\cL_0\) be lisse and
punctually pure of weight \(u\).  Then
\[
  \cL_0[d]\in\Perv(U_0)
  \quad\text{and is pure of weight }u+d.
\tag{10.4}\label{eq:lisse-perverse-weight}
\]
\end{proposition}

\begin{proof}
Perversity is \cref{prop:lisse-shift-perverse}.  The only ordinary
cohomology sheaf of \(\cL[d]\) is
\(\cH^{-d}(\cL[d])=\cL\).  For a complex of weight \(u+d\), the required
weight in ordinary degree \(-d\) is \((u+d)-d=u\), giving the upper bound.
Smooth duality gives
\[
  \D_U(\cL[d])\simeq\cL^\vee(d)[d].
\]
The dual local system has weight \(-u\), and the twist has weight \(-2d\);
therefore the right side has perverse-normalized weight \(-u-d\).  This is the
lower bound by \eqref{eq:derived-lower-dual}.
\end{proof}

In particular, on a smooth \(d\)-fold
\[
  \Qbarl[d]\quad\text{is perverse and pure of weight }d.
\tag{10.5}\label{eq:constant-perverse-weight}
\]
This single shift is the source of the final identity
\(d+(i-d)=i\).

\section{Ordinary versus perverse cohomology}

For a perverse sheaf \(\cP_0\), say that it has weights at most \(w\) if it
does so as a complex in the sense of \cref{def:derived-weight-bounds};
define lower bounds and purity similarly.

\begin{theorem}[BBD perverse weight criterion]
\label{thm:perverse-weight-criterion}
For a mixed complex \(\cK_0\),
\[
 \cK_0\in D_{\leq w}
 \quad\Longleftrightarrow\quad
 \pH^a(\cK_0)\text{ has weights }\leq w+a
 \text{ for every }a.
\tag{10.6}\label{eq:perverse-upper-equivalence}
\]
The same equivalence holds for lower bounds.  Consequently \(\cK_0\) is pure
of weight \(w\) if and only if
\[
  \pH^a(\cK_0)\quad\text{is pure of weight }w+a
\tag{10.7}\label{eq:pH-pure-shift}
\]
for every \(a\).
\end{theorem}

\begin{proof}
Filter \(\cK_0\) by its perverse truncation tower.  On a smooth stratum of
dimension \(d\), each perverse constituent is a shifted lisse sheaf
\(\cL[d]\), and \cref{prop:lisse-perverse-weight-shift} converts the desired
perverse inequality to the ordinary one.  For the closed complement, use
Noetherian induction and the recollement triangles.  Upper weight bounds pass
through extensions, subobjects, and quotients; duality gives the lower
statement.  The finite tower gives one implication by devissage and the other
by applying perverse cohomology to ordinary truncation triangles.  This is BBD
5.4.1; \eqref{eq:pH-pure-shift} is BBD 5.4.4.
\end{proof}

On the point \(\Spec\Fq\), perverse and ordinary cohomology coincide, so
\eqref{eq:pH-pure-shift} reduces to the familiar rule
\[
  \cK_0\text{ pure of weight }w
  \Longrightarrow H^a(\cK)\text{ pure of weight }w+a.
\tag{10.8}\label{eq:point-weight-shift}
\]

\section{Safe formal consequences}

\begin{lemma}[Weight bounds and finite filtrations]
\label{lem:weight-filtration-subquotients}
Let \(V\) carry geometric Frobenius and a finite Frobenius-stable filtration.
If every graded piece has weights at most \(w\), then \(V\) has weights at
most \(w\).  The analogous statement holds for lower bounds.  Hence every
subquotient of a mixed perverse sheaf of weights at most \(w\) has weights at
most \(w\).
\end{lemma}

\begin{proof}
Characteristic polynomials multiply in a short exact sequence, so the
multiset of eigenvalues is the union of those of the subobject and quotient.
Iterate through the filtration.  Costalks and stalks are exact on the
appropriate perverse ranges; the general assertion follows by the
stratification devissage used in \cref{thm:perverse-weight-criterion}.
\end{proof}

\begin{corollary}[Spectral-sequence rule]
\label{cor:perverse-weight-spectral-rule}
In a bounded Frobenius-equivariant spectral sequence, if every \(E_2^{a,b}\)
has weights at most \(u_{a,b}\), then every graded piece of the abutment has
the maximum of the bounds of the terms that can contribute to it.  No
degeneration is required.
\end{corollary}

\section{The dependency firewall}

The following important statements are \emph{not} available yet:
\begin{align*}
 \cP\text{ pure }w
 &\Longrightarrow j_{!*}\cP\text{ pure }w,\\
 \cK\in D_{\leq w}
 &\Longrightarrow Rf_!\cK\in D_{\leq w},\\
 \cK\text{ pure},\ f\text{ proper}
 &\Longrightarrow Rf_*\cK\text{ pure},\\
 \cK\text{ pure}
 &\Longrightarrow
   \cK_{\Fbar}\simeq\bigoplus_a\pH^a(\cK_{\Fbar})[-a].
\end{align*}
The second line is the main Deligne--Laumon theorem proved by the next four
chapters.  The first and third then follow by combining its upper bound with
duality; the fourth follows from weight separation.  Using any of these four
lines in the Fourier or curve argument would be circular.

\begin{perversebox}
Before the curve theorem, the allowed perverse tools are the \(t\)-structure,
recollement, \(j_{!*}\) as a minimal extension, Artin vanishing, Fourier
\(t\)-exactness, and the bookkeeping formulas
\eqref{eq:lisse-perverse-weight}--\eqref{eq:pH-pure-shift}.  No proper-purity
or decomposition statement is available.
\end{perversebox}

\begin{formalization}
Represent the three notions separately:
\begin{enumerate}
\item punctual weight of a constructible Weil sheaf;
\item upper/lower derived weight predicates;
\item perverse-heart weight predicates transported by the inclusion of the
      heart into the derived category.
\end{enumerate}
Then formalize \cref{thm:perverse-weight-criterion} as two equivalences and
\eqref{eq:pH-pure-shift} as their intersection.  Tests should verify that
shift \([1]\) raises weight by \(1\), twist \((1)\) lowers it by \(2\), and
\(\Qbarl[d]\) has perverse weight \(d\), not \(0\).
\end{formalization}
